\hypertarget{the-typing-system-static-analysis-vs.-runtime-enforcement}{%
\chapter{The Typing System: Static Analysis vs.~Runtime
Enforcement}\label{the-typing-system-static-analysis-vs.-runtime-enforcement}}

Python's type system has evolved from simple comments to a sophisticated
language-level feature. This chapter deconstructs how types exist in the
runtime.

\hypertarget{typing-internals-the-genericalias-and-specialform}{%
\subsection{\texorpdfstring{67.1 \texttt{typing} Internals: The
\texttt{GenericAlias} and
\texttt{SpecialForm}}{67.1 typing Internals: The GenericAlias and SpecialForm}}\label{typing-internals-the-genericalias-and-specialform}}

When you write \passthrough{\lstinline!list[int]!}, you are creating a
\passthrough{\lstinline!types.GenericAlias!} object. * \textbf{The
\passthrough{\lstinline!\_\_getitem\_\_!} Hook}: Classes like
\passthrough{\lstinline!list!} or \passthrough{\lstinline!dict!}
implement \passthrough{\lstinline!\_\_class\_getitem\_\_!} to support
the bracket syntax. * \textbf{Runtime Overhead}: Type hints are
evaluated at import time. Large-scale use of complex nested types can
noticeably slow down the startup of a Python application.

\hypertarget{static-vs.-runtime-verification}{%
\subsection{67.2 Static vs.~Runtime
Verification}\label{static-vs.-runtime-verification}}

\begin{itemize}
\tightlist
\item
  \textbf{Static Analysis}: Tools like \passthrough{\lstinline!mypy!} or
  \passthrough{\lstinline!pyright!} scan the AST (Chapter 31) and verify
  types without running the code.
\item
  \textbf{Runtime Enforcement}: Libraries like
  \passthrough{\lstinline!pydantic!} or
  \passthrough{\lstinline!beartype!} intercept function calls or class
  instantiation to verify types at execution time.
\item
  \textbf{\passthrough{\lstinline!inspect.get\_type\_hints()!}}: This
  function is the ``Godhood'' way to retrieve types at runtime, handling
  forward references (strings like \passthrough{\lstinline!"MyClass"!})
  by evaluating them in the correct namespace.
\end{itemize}

\hypertarget{protocols-and-structural-subtyping-pep-544}{%
\subsection{67.3 Protocols and Structural Subtyping (PEP
544)}\label{protocols-and-structural-subtyping-pep-544}}

Protocols allow for ``static duck typing.'' * \textbf{Internals}: A
\passthrough{\lstinline!Protocol!} class uses a specialized metaclass
that identifies which methods define the interface. Unlike
\passthrough{\lstinline!abc.ABC!}, you don't need to inherit from the
Protocol; you just need to implement the methods.

\begin{center}\rule{0.5\linewidth}{0.5pt}\end{center}
